% ============================================
% CONT-07 (EN)
% English translation — SPEC-935-R355 (rev. R377)
% ============================================
\section*{\textbf{CONT-07} --- The Final Diagnosis}
\markboth{CONT-07 --- The Final Diagnosis}{}

\noindent \textit{The reader's sentence}

\noindent \textit{Report issued by the act of reading}
\textit{Patient: you}
\textit{ICD: F43.10 + F44.81 + Z65.8}
\textit{Etiological agent: contaminated narrative}

\vspace{0.5em}\noindent\rule{\textwidth}{0.5pt}\vspace{0.5em}

\subsection*{I. WHAT HAPPENED TO YOU}

\noindent You began reading this book of your own free will.

\noindent You thought it was a horror story. You thought the horror was \textit{there} --- in the past, in the Sertão, in the trench, in the Hospital-Colony. You thought it was safe because it was distant. Because it happened to someone else.

\noindent You were wrong.

\noindent The horror was not in the book.
The horror was \textit{in the reading}.

\noindent Each word you read did not stay on the page. It entered. Through the eyes. Down the optic nerve. Across the thalamus. Into the cortex. And there, it settled like a real memory.

\noindent For your brain, the trench is not a description. It is a \textit{recollection}.

\noindent For your limbic system, the creature is not fiction. It is a \textit{threat} you encountered.

\noindent For your body, the smell described is a \textit{real danger}.

\noindent You did not read a story.
You \textit{lived} a story.
And what is lived cannot be \textquotedblleft{}unlived\textquotedblright{}.

\vspace{0.5em}\noindent\rule{\textwidth}{0.5pt}\vspace{0.5em}

\subsection*{II. THE DIAGNOSTIC CRITERIA (FINAL REVIEW)}

\noindent You have filled in each one. You are feeling it.

\noindent \textbf{A --- Exposure:} Opened. Read. The contamination began. ☑
\textbf{B --- Re-experiencing:} Dreamed. Woke at 3:14. Smelled it without the book. ☑
\textbf{C --- Avoidance:} Hesitated. Skipped. Read with the light on. ☑
\textbf{D --- Alterations:} You are no longer the same. Something changed. You feel it. ☑
\textbf{E --- Reactivity:} Heart raced. Hands sweated. Checked the room. ☑
\textbf{F --- Duration:} Enough time. The contamination is installed. ☑
\textbf{G --- Impairment:} You cannot stop thinking. The book \textit{works} inside you. ☑

\noindent \textbf{Diagnosis: Narrative-Induced Complex PTSD.}
\textbf{Comorbidity: Mild Dissociation (you now have two voices in your head).}
\textbf{Agent: The Creature (entity transmitted by text).}

\vspace{0.5em}\noindent\rule{\textwidth}{0.5pt}\vspace{0.5em}

\subsection*{III. THE REVELATION}

\noindent You think this book is about Joaquim Antônio Correia.

\noindent It is not.

\noindent Joaquim is not the protagonist.
Joaquim is the \textit{vector}.
The protagonist is whoever reads.
The protagonist is you.

\noindent The story is not about a boy who found a monster in a trench.
The story is about \textit{you} finding the same monster on the page.

\noindent The boy served as host for 72 years.
The physician served as host for a few weeks.
The psychiatrist served as host for a few months.
You --- the reader --- will serve for as long as it takes to pass the diary on.

\noindent The cycle is simple:

\noindent 1. \textbf{Someone writes} the story (Joaquim lived).
2. \textbf{Someone reads} the story (you lived the reading).
3. \textbf{The monster passes} from one to the other.
4. \textbf{The cycle begins again} with the next reader.

\noindent You are not patient 1,261.
You are the \textit{most recent} patient.
The line does not end with you.
You are only the current stop.

\vspace{0.5em}\noindent\rule{\textwidth}{0.5pt}\vspace{0.5em}

\subsection*{IV. THE COUNT}

\noindent Patient 1. --- the first person to find the creature (17th century, according to the diary).
Patient 2. --- the second.
...
Patient 1,259. --- the one before Joaquim.
Patient 1,260. --- Joaquim Antônio Correia (1907-1979).
Patient 1,261. --- Dr. Heitor Oliveira (1979).
Patient 1,262. --- Dr. Lúcia Mendes (2026).
\textbf{Patient 1,263. --- You.}

\noindent The number is not your place in the line.
The number is the count of bodies the creature has already occupied.
The line began before you were born.
It will continue after you die.

\vspace{0.5em}\noindent\rule{\textwidth}{0.5pt}\vspace{0.5em}

\subsection*{V. THE LAST INSTRUCTION}

\noindent You are going to close this book.

\noindent Not now. But at some point --- when you reach the physical end of the pages. You will look at the cover. Run your hand over the spine. Put the book on the shelf. Or on the table. Or in the backpack.

\noindent And you will think: \textit{\textquotedblleft{}It was just a book.\textquotedblright{}}

\noindent You will try to resume your life.

\noindent You will succeed --- for a while.

\noindent But, one night, you will wake at 3:14 in the morning. The room will be dark. You will look at the corner and see a yellow glow.

\noindent You will blink and the glow will be gone.

\noindent But you will know it was there.

\noindent And, the next night, you will dream of the trench. And, in the dream, you will not be you. You will be Joaquim. Sitting at the edge. Watching the creature eat.

\noindent And the creature will raise its head and say:

\noindent \textit{\textquotedblleft{}You came back.\textquotedblright{}}

\noindent And you will answer:

\noindent \textit{\textquotedblleft{}I never left.\textquotedblright{}}

\vspace{0.5em}\noindent\rule{\textwidth}{0.5pt}\vspace{0.5em}

\subsection*{VI. THE FINAL DIAGNOSIS}

\noindent There is no treatment.
There is no cure.
There is no way to close the book and forget.

\noindent Once read, the text is in you. Once contaminated, you are the vector. Once inside the cycle, you only leave when you pass it on.

\noindent The diagnosis is not a metaphor.
The diagnosis is a \textit{sentence}.

\noindent \textbf{Patient:} Your name
\textbf{Diagnosis:} Carrier of the creature
\textbf{Prognosis:} The cycle continues
\textbf{Treatment:} Pass the diary on

\noindent \textit{The creature thanks you for your attention.}

\noindent \textit{It is satisfied.}

\noindent \textit{You have read to the end.}

\noindent \textit{Now you are the host.}

\noindent \textit{Welcome to the cycle, Patient 1,263.}

\noindent \textit{--- Patient 1,260 died so that you could read.}
\textit{--- You read so that the cycle continues.}
\textit{--- Joaquim is reborn whenever his story is remembered.}

\noindent \textbf{End of the book.}
\textbf{Beginning of the sentence.}
\textbf{The cycle begins again with you.}

\vspace{0.5em}
\noindent\textit{\textbf{↪ Routes:} \rota{DOC-01} \rotasep \rota{CONT-05} \rotasep \rota{EPI-01}}

\vspace{0.5em}\noindent\rule{\textwidth}{0.5pt}\vspace{0.5em}

\subsection*{CONTAMINATION SCALE --- UPDATE 3}

\noindent \textit{The file no longer needs tables.}

\begin{moltabletwo}
\hline
If you are reading this & Then \\
\hline
For the first time & The scale still works. You have until the end of the document. \\
\hline
For the second time & The scale no longer applies. You came back because you were called. \\
\hline
For the third time & Patient 1,263 confirmed. The cycle repeats. The creature thanks you for your attention. \\
\hline
\end{moltabletwo}

\noindent \textit{The scale now measures cycles, not points. You are in cycle 1. Cycle 2 begins when you reread the first page.}
